% ======================================================================
%  53b_dialogue_B.tex — 附录 D 第二部分《线算子与特征膜》问答录
%  底本：KW 2006（hep-th/0604151），§6--§8，印刷 pp.78--122
%  只收原文明确陈述的 A 级内容；每条答案可回原文核对，短引 <=120 字符
% ======================================================================
\subsection{线算子与特征膜：\S6--\S8（pp.78--122）}

本录只收原文 \S6--\S8 明写的内容：线算子与圈算子的定义及其物理意义、紧化到二维后它们变成什么、\S7 的离散通量与 $S$-对偶、以及 \S8 中 Wilson 算子如何作用于膜。问题按原文页序排列，答案紧贴原文措辞；凡属我们的推测、或原文只是暗示而非明说的，一律不收。每条末尾的英文引文逐字取自原文（必要时截成语段），方括号中标出它在印刷本中的页码，可据以回查底本。

\QAq{\S6 为什么把算子挂到一维流形上？}
\QAa{原文说这套四维拓扑场论的特殊之处，在于它把算子关联到 $M$ 中定向的一维流形 $S$ 上，而算子只依赖 $S$ 的同伦类、不只是同调类。作者随即点明动机：重点并不是研究 $M$ 的基本群，而是这些一维流形算子与膜之间表现出很有意思的关系。}
\QAs{src p.78}{``these one-manifold operators prove to have very interesting properties in relation to branes''}

\QAq{什么情形下叫圈算子，什么情形下才有线算子？}
\QAa{原文规定：$M$ 紧致且无边界时，$S$ 必须是闭圈，这类一维流形算子因此通称圈算子。若四维流形非紧，则允许延伸到无穷远的线算子；其中最要紧的是在某个空间点贯穿全部时间的静态线算子，它在量子化时必须被包括进来。}
\QAs{src p.79}{``If $M$ is compact and without boundary, then $S$ is required to be a closed loop''}

\QAq{Wilson 圈算子的定义式是什么？}
\QAa{原文取 $G$ 的不可约表示 $R$ 与 $M$ 中定向的圈 $S$，定义 $W_0(R,S)=\Tr_R P\exp(-\oint_S A)$，即 $A$ 绕 $S$ 的和乐在 $R$ 中的迹。负号来自该书的约定 $D=\dd+A$；若把规范场写成 hermitian 的 $A'=\ii A$，指数就变成 $+\ii\oint A'$。要在扭化理论里保持拓扑对称，须取 $t=\pm\ii$，把 $A$ 换成 $A+\ii\varphi$ 或 $A-\ii\varphi$：此时拓扑方程恰好断言它平坦，其和乐才是自然的不变量。}
\QAs{src p.80}{``$W_0(R,S)$ is the trace, in the representation $R$, of the holonomy of $A$ around $S$''}

\QAq{开曲线上的 Wilson 算子为什么不取迹？}
\QAa{原文给出开线情形的定义：若 $S$ 有端点 $p$、$q$，就把 $W(R,S)$ 定义为平行输运矩阵而非迹，它是 $p$ 处纤维到 $q$ 处纤维的映射；迹只在闭圈情形出现。作者特别指出，这个推广版本在真正进入几何 Langlands 时是不可缺的。}
\QAs{src p.80}{``we define $W(R,S)$ not as a trace but as the matrix of parallel transport''}

\QAq{'t Hooft 算子与 Wilson 算子是什么关系，它怎样定义？}
\QAa{原文把 Wilson 算子称为 order 算子，并说 order 算子的对偶通常是 disorder 算子；在本问题里，Wilson 算子的对偶正是 't~Hooft 算子。在 abelian 的基本情形里，'t~Hooft 算子由路径积分定义：只对 $M\setminus S$ 上方线丛的联络积分，并要求它们在 $S$ 附近与奇异的 Dirac 单极子相符，使曲率在 $S$ 附近按 $F\sim(\ii/2)\star_3\dd(1/r)$ 发散——这就是“用奇异规范场定义”的含义。与 Wilson 需要 $S$ 的定向不同，'t~Hooft 圈算子需要的是 $S$ 法丛的定向；荷 $m$ 的算子就把基本奇点取 $m$ 倍。}
\QAs{src p.82}{``the path integral is to be performed over connections on a line bundle over $R^4\setminus L$''}

\QAq{为什么 't Hooft 奇点必须让 $\varphi$ 也奇异？}
\QAa{因为扭化理论的拓扑方程 $(F+\dd\varphi)^+=(F-\dd\varphi)^-=0$ 在 $F$ 有 Dirac 单极奇点而 $\varphi$ 光滑时无解。原文因此要求场在 $S$ 附近取 (6.9) 的奇点：$F=(\ii/2)\star_3\dd(1/|x|)$、$\varphi=(\ii/2|x|)\dd s$，其中 $s$ 是沿 $S$ 的“时间”坐标；荷 $m$ 即取 $m$ 倍。}
\QAs{src p.83}{``These equations are not satisfied if $F$ has the Dirac monopole singularity and $\varphi$ is non-singular''}

\QAq{任意紧群 $G$ 的 't Hooft 算子怎样分类？}
\QAa{原文任取一个同态 $\rho:U(1)\to G$，要求场在 $S$ 附近的奇点是 (6.9) 在 $\rho$ 下的像；算子只依赖 $\rho$ 的共轭类。共轭类同态由对偶群 $\Ldual$ 的最高权分类，所以每个 $\Ldual$ 的不可约表示 $\LR$ 给出一个 't~Hooft 算子 $T(\LR,S)$，这就是经典的 GNO 对偶。同一条对应还使 't~Hooft 算子的分类等同于 $\PP^1$ 上 $G$-丛的分类。}
\QAs{src p.84}{``The fact that 't~Hooft operators of $G$ are classified by representations of $\Ldual$ is the classic GNO duality''}

\QAq{紧化到二维后，四维的圈算子变成什么？}
\QAa{原文取 $M=\Sigma\times C$（$\Sigma$ 远大于 $C$，先设 $\Sigma$ 无边）。若圈形如 $S=p\times S'$（$p\in\Sigma$、$S'\subset C$），它在 $\Sigma$ 上的有效理论里就是一个点状局域算子；若 $S=S''\times q$，它仍是二维的圈算子；一般的曲线同伦于在 $\Sigma$ 与 $C$ 之间交替传播的曲线，投影到 $\Sigma$ 上便成为带局域算子插入的圈算子。}
\QAs{src p.90}{``simply looks like a pointlike local operator in the effective theory on $\Sigma$''}

\QAq{二维里线算子为什么可以用膜来理解？}
\QAa{原文指出，二维流形被一维流形局部地分成两个不相交的区域，线算子两侧的耦合可以不同，因此它可能产生长程效应。用折叠技巧把两侧放到同一侧后，沿算子折叠出的边界条件就是一个 $(Y,Z)$-膜；而这里所说的“膜”，一般就是指量子场论里的局域边界条件。于是线算子在二维里归结为一个膜。}
\QAs{src p.91}{``the line operator reduces in two dimensions to a $(Y,Z)$-brane''}

\QAq{线算子靠近边界时起什么作用？}
\QAa{原文把线算子 $X$ 与膜 $B$ 合成的有效边界条件记作 $B'=XB$，即 $X$ 作用在旧膜上产生一个新膜。因此可以把线算子看成作用在边界条件上、产生新边界条件的对象。这正是几何 Langlands 关注的情形：$\Sigma$ 有边界，线算子与边界平行。}
\QAs{src p.94}{``we write $B'=XB$ for the new brane that is produced by acting with $X$ on the old brane''}

\QAq{什么叫线算子的特征膜（eigenbrane）？}
\QAa{原文的定义是：若对某个向量空间 $V$ 有 $XB=B\otimes V$，就称 $B$ 是线算子 $X$ 的特征膜。做法是：对每个膜 $B$ 与向量空间 $V$ 都定义一个新膜 $B\otimes V$；本征值应看成一个向量空间而不是数（$V=\CC^k$ 时只是取 $k$ 份）。原文说这个概念正对应 Beilinson--Drinfeld 在几何 Langlands 中定义的特征层。}
\QAs{src p.95}{``A brane $B$ is an eigenbrane for a line operator $X$ if $XB=B\otimes V$, for some vector space $V$''}

\QAq{电特征膜、磁特征膜是什么？$S$-对偶怎样互换二者？}
\QAa{原文定义：Wilson 线算子的联合特征膜叫电特征膜，'t~Hooft 算子的联合特征膜叫磁特征膜。因为四维来的线算子可在 $C$ 上任取一点 $p$，它们自动对易；而 $V_p$ 局部不依赖 $p$、只有整体单值，故拼成 $C$ 上的平坦向量丛——所以电、磁特征膜都自动带着 $C$ 上的平坦向量丛。$S$-变换 $S:\Psi\mapsto-1/(n_g\Psi)$ 把拓扑 Wilson 算子送到拓扑 't~Hooft 算子，因而把电特征膜送到磁特征膜。}
\QAs{src p.97}{``A joint eigenbrane of the Wilson line operators will be called an electric eigenbrane''}

\QAq{离散电通量与离散磁通量是怎样定义的？}
\QAa{原文先指出中心 $Z=\pi_1(G_{\mathrm{ad}})$ 自然自对偶，存在 perfect pairing $\Upsilon:Z\times Z\to U(1)$。在 $M=S^1\times W$ 上，丛的拓扑由 $a\in H^1(W,Z)$ 与 $m\in H^2(W,Z)$ 给出；character $e:H^1(W,Z)\to U(1)$ 与 $m$ 就分别称为离散电通量与离散磁通量，而 $\sum_a e(a)Z_{a,m}$ 可在 Hilbert 空间 $\mathcal H_{e,m}$ 里读成迹。$S$-对偶交换 $E$ 与 $M$，把 $(e,m)$ 送到 $(-m,e)$。}
\QAs{src p.105}{``$e$ and $m$ are called respectively the discrete electric and magnetic flux''}

\QAq{什么叫 universal bundle？它一定存在吗？}
\QAa{原文定义：universal bundle 是 $G$-丛 $\mathcal E\to M\times C$，使每个 $p\in M$ 上 $\mathcal E|_{p\times C}$ 都同构于 $E_p$。对伴随群 $G_{\mathrm{ad}}$ 粘合没有障碍，故 $\mathcal E_{\mathrm{ad}}$ 存在；但对一般 $G$，$\Theta_{ki}\Theta_{jk}\Theta_{ij}=f_{ijk}\in Z$ 给出非平凡元素 $\zeta$，它是 universal $G$-丛存在的阻碍。原文补充：它总可以作为 twisted vector bundle 存在，twist 为 $\varrho(\zeta)$。}
\QAs{src p.109}{``it does always exist as a twisted vector bundle, twisted by $\varrho(\zeta)$''}

\QAq{\S7.2 里 $m_0$ 与 $e_0$ 各是什么？}
\QAa{原文说靶空间 $\MH(G_{\mathrm{ad}},C)$ 并不连通，其连通分支由 $G_{\mathrm{ad}}$-丛 $E\to C$ 的拓扑型标记，$m_0$ 就是这个拓扑型。$e_0$ 则给出 $\MH(G_{\mathrm{ad}},C)$ 上的平坦 $B$-场 $b_{e_0}=e_0(\zeta)$，它在有效 sigma 模型里的作用正是按平坦 $B$-场的样式给每个映射加权；$e_0\ne0$ 时，膜的 Chan-Paton 丛是 twisted vector bundle。}
\QAs{src p.111}{``Its components are labeled by the topological type of the $G_{\mathrm{ad}}$-bundle $E\to C$''}

\QAq{\S7.2 后半的 $e_1$、$m_1$ 度量什么？}
\QAa{原文把 $\Sigma=S^1\times\tilde S^1$ 时 $\xi(E)$ 的其余分量记作 $a_1$、$m_1$：$m_1$ 度量 $\Phi|_{\tilde S^1}$ 在基本群中的同伦类，$a_1$ 度量沿时间方向的绕数。在二维图像里，$\MH(G)$ 上的弦有守恒动量 $E_C=H^1(C,\ZZ)$，$\MH(G_{\mathrm{ad}})$ 上的弦有守恒绕数 $M_C=\pi_1(\MH(G_{\mathrm{ad}}))=H^1(C,\ZZ)$。$S$-对偶交换这两套守恒量，原文说这正是它与二维 $T$-对偶有关的体现。}
\QAs{src p.115}{``S-duality or Montonen-Olive or Langlands duality exchanges the two pictures''}

\QAq{\S8.1 里 Wilson 线怎样作用于膜？}
\QAa{原文结论：表示 $R$、支在点 $p\in C$ 的 Wilson 线把膜 $B$ 的 Chan-Paton 丛 $U$ 变成 $U\otimes E_p(R)$。$E_p(R)$ 一般不是向量丛，而是 twist 为 $\theta_R(\zeta)$ 的 twisted vector bundle，所以这个作用把 $B$-场改成 $b\to b+\theta_R(\zeta)$，也就是把 \S7 的离散电场改成 $e_0\to e_0+\theta_R$。$S$-对偶的像则是 't~Hooft 线把 $m_0$ 改成 $m_0+\xi(\LR)$。}
\QAs{src p.119}{``it changes the discrete electric field studied in section 7 by $e_0\to e_0+\theta_R$''}

\QAq{为什么支在一点的零膜是电特征膜？它与 $\pi_1(C)$ 的表示有什么关系？}
\QAa{原文论证：任何向量丛限制到一点都平凡，故 $E_q(R)$ 限制到 $x\in\MH$ 也平凡，于是 $U_x\otimes E_q(R)=U_x\otimes E_q(R)|_x$；零膜 $B_x$ 就是电特征膜，本征值是向量空间 $E_q(R)|_x$。在几何 Langlands 一侧，起手是 $\Ldual$ 与同态 $\vartheta:\pi_1(C)\to\Ldual_\CC$，它定义 $\MH(\Ldual,C)$ 中的点 $x(\vartheta)$；若 $\vartheta$ 不可约，则 $x(\vartheta)$ 是光滑点，支在它上面的零膜就是电特征膜。原文还指出，电特征膜其实是不同 $e_0$ 的 twisted 零膜之和。}
\QAs{src p.121}{``of the Wilson loop operator $W_q(R)$ being the vector space $E_q(R)|_x$''}
